%TODO 
%\section{Invariances}
%Section about invariances and why they are needed in localization.

%The "filter" must give the same response invariant of the position in the image 
%and its often wanted to have it rotational and scale invariant (if your not 
%intending to measure things like angles or size that is)

\section{Data}
\subsection{Image}
In a computer an \emph{image} is represented as an array of numbers, mostly 8bit integers ranging between $\left[0,255\right]$ 

\subsubsection{RGB image}
A color image can be defined as
\begin{equation}
    \begin{array}{ccc}
    RGBimage : \NN^2 &\rightarrow& \NN^3\\
    RGBimage : position&\rightarrow&color
    \end{array}
\end{equation}
Which basically is foreach position in the image we have a response color.
An alternitve represenation whould be $\NN^5=\langle pos,value\rangle$.

The RGB image space is denoted by
\begin{equation}
    RGBimage\in \IS_{RGB}
\end{equation}

\subsubsection{Grayscale image}
A grayscale image can, similary to RGB images, be defined as a function
\begin{equation}
    \begin{array}{ccc}
    I : \NN^2 &\rightarrow& \NN\\
    I : position&\rightarrow&intensity
    \end{array}
\end{equation}
or alternatively $\NN^3=\langle pos,intensity\rangle$.

Finally the $image$-space is denoted by
\begin{equation}
    I \in \IS
\end{equation}

\begin{example}~\\
    \begin{tabular}{lcr}
        $
        \begin{pmatrix} 
        255&  255&  255&  255&  255&  255&  255&  255&  255&  255\\
        255&  255&  255&  255&  255&  255&  255&  255&  255&  255\\
        255&  255&  255&    0&  255&  255&    0&  255&  255&  255\\
        255&  255&  255&    0&  255&  255&    0&  255&  255&  255\\
        255&  255&  255&  255&  255&  255&  255&  255&  255&  255\\
        255&    0&  255&  255&  255&  255&  255&  255&    0&  255\\
        255&  255&    0&  255&  255&  255&  255&    0&  255&  255\\
        255&  255&  255&    0&    0&    0&    0&  255&  255&  255\\
        255&  255&  255&  255&  255&  255&  255&  255&  255&  255\\
        255&  255&  255&  255&  255&  255&  255&  255&  255&  255
        \end{pmatrix}
       $
       &$\Rightarrow$& \parbox[c]{1em}{\includegraphics[scale=10.0]{smiley.png}}
    \end{tabular}
\end{example}

\subsubsection{Video}
A \emph{video} can be represented as a function mapping an integer to the designated image $I$
\begin{equation}
    video : \NN \rightarrow I
\end{equation}

\subsection{State}
<intro about state>
\begin{note}
    State and hypothesis is model dependant.
\end{note}
\subsubsection{State}
$x$
\subsubsection{State-space}
$x\in\XS$ %TODO(collides with X?)
\subsubsection{Hypothesis}
$z$

Explain: Hypotesis (in the context of PF), sample, DOF?, 

\subsubsection{Hypothesis-space}
$z\in\ZS$

\subsubsection{Estimate}
The algorithms estimate of the ground truth
$x^*$ 

\subsubsection{Degrees of Freedom (DOF}
The number of adjustable parameters in a model.

\subsubsection{Observation}
<placement?><denotation?>

\subsection{MISC?}

\subsubsection{Prop. function as a collection of points with attached weight.}

\subsubsection{Dynamic system}
    It's a tuple <space,update,time>
    Any mechanical system following newton physics (or relativistic for that matter) can be considered to be an dynamic system
    in our case the tuple whould be <feature\_space,update\_rule (the one we are trying to preprocess data to approximate),time>

